% GENERATED by scripts/md2latex_manuscript.py -- do not edit by hand.
% source: papers/maskfree_bundle/manuscript.md

\section{Introduction}


Dynamic content breaks monocular and RGB-D SLAM in a specific way: a moving object supplies photometric and geometric residuals that the estimator, which assumes a rigid world, resolves by moving the camera. On the standard dynamic benchmarks this is not a marginal degradation. Our 3DGS-SLAM backbone, run unmodified on 3 seeds, drifts to 39.32$\pm$1.01 cm on a sequence with a person and a balloon, 43.85$\pm$8.55 cm on a person-tracking sequence, and 147.5$\pm$39.0 cm on a crowded one --- against roughly 1.5 cm on the static sequences from the same datasets.


The dominant response is to detect the movers and remove them. An external, pretrained, class-conditioned network labels which pixels belong to things that could move, and the estimator discards them; this structure runs from keypoint-era systems through neural-implicit SLAM and into essentially every recent 3DGS-based dynamic system (Sec. 2.1). It works, and we do not argue otherwise. But it makes the dynamic decision a per-frame classification by a second model, with that model's runtime cost and failure modes, and it is closed-set: the category list has no label for many of the movers these benchmarks actually contain --- balloons, carried boxes, objects placed and removed. A recent line replaces the \emph{class} prior with a learned uncertainty (Sec. 2.2), and a third keeps a segmentation network but moves the dynamic decision out of it (Sec. 2.3).


Our own configuration sits in that third line, and we want to be exact about what is and is not new in it. It decides which pixels to trust from \emph{consistency} --- the disagreement between observed optical flow and the flow a rigid scene would produce --- rather than from categories, and it down-weights rather than removes. \textbf{Operating without semantic segmentation is not itself a contribution we claim}; prior work has done that (Sec. 2.3). What has not been supplied, for this family of designs, is evidence answering two questions that decide whether such a system is understood at all. \emph{Which} of its parts are actually necessary, and on which sequences do they stop being necessary --- a system reported as an assembled bundle leaves this unmeasured. And is a reported gain attributable to the proposed mechanism, or merely to the extra information the mechanism was handed --- which can only be separated by a counterpart consuming \emph{the same information} through a deliberately naive mechanism.


We therefore present a mask-free kernel for dynamic 3DGS-SLAM together with the measurements that bound it, and our contributions are the measurements:


\begin{enumerate}
\item \textbf{Component necessity is regime-dependent, and we measured it rather than asserting it.} A pre-registered three-factor factorial (8 cells $\times$ 5 sequences $\times$ 3 seeds = 120 runs) finds that uniform keyframe coverage is strongly positive on object-only sequences but \textbf{negative} on the mixed-mover sequence; that robust reweighting is indistinguishable from zero on 3 of 5 sequences yet is the single largest effect in the matrix on the hardest one; and that the reliability signal is the only component never negative anywhere. \textbf{No sequence has all three positive --- 0 of 5 --- so no globally necessary subset exists} in this design (Sec. 5.2). This is a statement about mechanism, not engineering: the necessary subset is a function of the regime.

\item \textbf{The gain is not attributable to the information alone.} Under an equal, byte-identical, strictly causal flow budget, a naive per-pixel threshold on the \emph{same} frozen flow improves on the baseline on only 2 of 4 held-out sequences and is \textbf{worse than the baseline on one} (0.90$\times$), while our mechanism improves on 4 of 4 (1.63--4.99$\times$) and beats the naive counterpart on 4 of 4 (Sec. 5.5). This closes the ``it just had more information'' reading.

\item \textbf{The pose-side mechanism is a variance--bias trade, confirmed by measuring the intermediate quantity.} One single-variable intervention moves per-frame tracking error \emph{worse} (+0.0806, 1.22$\times$ its noise floor) and full-trajectory error \emph{better} ($-$1.42 cm, 4.0$\times$), and the coherent, drift-contributing component of the error explains why by falling ($-$0.0248, 1.38$\times$) (Sec. 5.7). The consequence generalises beyond our system: \textbf{a claim that pose estimation improved is not well formed until the endpoint is named}, because a per-frame-magnitude endpoint does not predict trajectory quality.
\end{enumerate}


We also report where our mask-free configuration is \emph{not} sufficient, since the same evidence that supports the above delimits it. Semantic segmentation turns out to be redundant on both pure-object sequences and on one of two person sequences, but dominant on mixed-mover sequences and on a second person sequence where our mask-free configuration fails outright (32.41$\pm$8.51 cm, recovered to 11.89$\pm$2.36 cm with the mask). Across all 18 sequences the mask-free arm's own per-frame drift separates those two groups without overlap (Sec. 5.6), and its pre-registered midpoint held on 5 sequences unseen during development --- a diagnostic we report with its limits rather than as a selector. The regime split also survives the converse test: starting from a simpler configuration, adding coverage alone recovers the full configuration's accuracy on the object-only sequence, while on the crowd sequence no single component does (Sec. 5.3). We do not claim state-of-the-art accuracy: our absolute error trails the best reported dynamic systems by 1.4--2.5$\times$ on the TUM walking and person sequences, and we give the full 18-sequence table without selective highlighting.




\section{Related Work}


\subsection{Dynamic SLAM built on explicit dynamic priors}


Handling moving objects in SLAM has long followed a \emph{detect-then-exclude} structure: an external model labels which pixels belong to potentially moving things, and the estimator discards them. In the keypoint era DynaSLAM\textasciitilde{}\cite{dynaslam} computes ``a pixel-wise semantic segmentation of the images'' with Mask R-CNN\textasciitilde{}\cite{maskrcnn}, and NGD-SLAM\textasciitilde{}\cite{ngdslam} keeps the structure with a YOLO detector; neural-implicit SLAM inherited it, with RoDyn-SLAM\textasciitilde{}\cite{rodynslam} forming its motion mask by ``incrementally fusing the optical flow mask'' with ``the semantic segmentation mask'' to reject sampled rays.


The pattern is widespread in 3DGS-based dynamic SLAM. RGD-SLAM\textasciitilde{}\cite{rgdslam} segments with ``the pre-trained version of OneFormer on the COCO dataset'' and refines the motion mask with an EKF; DG-SLAM\textasciitilde{}\cite{dgslam} uses OneFormer ``to generate prior semantic segmentation'' and initialises poses from DROID-VO\textasciitilde{}\cite{droidslam}; BDGS-SLAM\textasciitilde{}\cite{bdgsslam} turns YOLOv5 detections into a per-Gaussian Bayesian dynamic prior; DAGS-SLAM\textasciitilde{}\cite{dagsslam} ``fuses lightweight YOLO instance priors'' with a per-Gaussian motion probability; and DynaGSLAM\textasciitilde{}\cite{dynagslam} combines ``coarse motion blobs computed from the RAFT\textasciitilde{}\cite{raft} optical flow image with real-time online SAM2'', adopting poses from DynoSAM ``without modification''. Every quotation here is anchored to its source line in the evidence appendix.


Whatever the representation, the dynamic decision here is delegated to an external, pretrained, class-conditioned model. That inherits a second model's runtime cost and failure modes; it makes the decision a \emph{per-frame classification}, so a missed or spurious detection propagates into the map; and it is \emph{closed-set}, with no label for many movers these benchmarks actually contain --- balloons, carried boxes, objects placed and removed.


\subsection{Replacing the class prior with a learned uncertainty}


A recent alternative drops the category prior while keeping a learned one. WildGS-SLAM\textasciitilde{}\cite{wildgsslam} predicts per-pixel uncertainty from ``a pre-trained DINOv2 feature extractor'' --- the finetuned version ``which injects 3D awareness'' --- decoded by a shallow MLP and supervised with metric depth ``as estimated by Metric3D v2''. Its own baseline overview defines \emph{Prior Free} as ``not using class priors, where 'O' represents object detection and 'S' denotes semantic segmentation'', which places the contribution precisely: the \emph{semantic class} prior is removed, while a learned appearance prior remains, supplied by two external pretrained networks.


This distinction matters for how we position our own system, and we hold ourselves to it: our mask-free configuration removes semantic segmentation, \textbf{not} the dependence on a pretrained network --- it still consumes frozen optical flow. We therefore describe it as \emph{mask-free}, never as \emph{prior-free}.


\subsection{Dynamic decisions from consistency rather than from categories}


A third line keeps a segmentation or depth network but moves the \emph{dynamic decision} out of it, or removes the semantic network entirely. Gassidy\textasciitilde{}\cite{gassidy} separates objects from background through ``the instance masks Si generated by YOLO segmentation'', then decides which of those objects are dynamic by iteratively analysing rendering-loss flows from a photometric-geometric loss, so the category label never decides motion. DGS-SLAM\textasciitilde{}\cite{dgsslamkong} runs an online Track-Anything segmentation in its frontend and adds a robust mask enforcing photometric consistency across keyframes, explicitly ``reducing noise from inaccurate segmentation and artifacts such as shadows''. Dy3DGS-SLAM\textasciitilde{}\cite{dy3dgsslam} drops the semantic network in a monocular setting, fusing an optical-flow mask with a depth mask from ``the DepthanythingV2 estimation network''. GARAD-SLAM\textasciitilde{}\cite{garadslam} goes furthest: it states that ``we do not perform semantic segmentation on the feature points in the tracking phase'', and instead builds a conditional random field over Gaussians in the backend, maps the segmentation back to frontend feature points, and validates labels with sparse optical flow.


Our own configuration belongs to this line, next to GARAD-SLAM. \textbf{We therefore do not claim that operating without semantic segmentation is by itself new.} What this line does show is that ``mask-free'' has to be checked per system rather than read off a label: of these four, two still consume a segmentation network and a third a monocular depth network.


\subsection{Robust estimation and temporal sampling in SLAM}


The ingredients we combine are individually classical. Robust M-estimation with Huber or Cauchy kernels is standard practice for suppressing gross outliers in photometric and geometric residuals; residual-based reweighting has ample precedent in dense and neural SLAM; and keyframe management is a core component of every modern system. Our backbone, MonoGS\textasciitilde{}\cite{monogs}, ``maintain[s] a small window W\_k'' of keyframes over which Gaussians and camera poses are optimised, states that for camera localisation it ``does not use any pretrained monocular depth predictor or other existing tracking modules'', and reports its accuracy as ``comparable to systems which perform explicit loop closure'' - i.e. it performs none itself.


What is missing is therefore neither another robust kernel nor the removal of the semantic prior, but the two measurements Sec. 1 named: which components are individually necessary, and whether a gain is attributable to the mechanism or to the information it was handed. Sec. 5 supplies both.


\subsection{Reading absolute numbers across these systems}


The systems above are not evaluated under a common protocol. They differ in input modality (WildGS-SLAM is monocular RGB, the rest RGB-D), in which pretrained networks they consume, in how far back later frames are allowed to rewrite the reported trajectory --- several run a global or dense bundle adjustment after tracking completes, whereas MonoGS and our system optimise only within a bounded keyframe window --- and in reported hardware, which one system does not state. We list these axes per method, with the supporting sentence and line number for each cell, in S6, and our own configurations appear there under the same columns. We do this to make the comparison legible, not to argue that any of these methods is unfair or that our numbers should be discounted; \textbf{no conclusion in this paper rests on that table.}




\section{Method}


Our system adds four components to a 3DGS-SLAM backbone. None of the four is a new primitive, and we do not present them as such (Sec. 2.4); this section states precisely \textbf{what each one changes in the estimator}, so that the ablations in Sec. 5 can be read as interventions on named quantities rather than on a monolithic ``our method''. Every constant quoted below is the resolved value actually used in the reported runs, not a nominal default (Sec. 3.7).


\subsection{Backbone and where the components attach}


The backbone is MonoGS\textasciitilde{}\cite{monogs} in its RGB-D configuration: the frontend estimates each frame's pose by minimising a photometric and geometric residual against a fixed Gaussian map, and the backend periodically re-optimises Gaussian parameters and the poses of a bounded keyframe window. We change neither the Gaussian representation nor the rasteriser. We \textbf{do} change the density lifecycle, and we say so rather than describing our configuration as the stock backbone.


Four components attach at four distinct places, named separately because Sec. 5.2 shows their contributions are not interchangeable (Fig. 1): \textbf{coverage} decides \emph{which frames become keyframes} (\S{}3.2); \textbf{robust reweighting} decides \emph{how much each pixel's residual counts in the pose solve} (\S{}3.3); \textbf{reliability weighting} decides \emph{which pixels are trusted to be static} (\S{}3.4); and the \textbf{lineage-tagged density lifecycle} decides \emph{which newly created Gaussians survive} (\S{}3.5). The third is the only one consumed at more than one site --- one in the frontend, one in the lifecycle --- and \S{}3.4 names both, because the ablation of Sec. 5.2 removes them together.


A fifth, \textbf{optional} component supplies a semantic person mask (\S{}3.6). Our \emph{mask-free} configuration disables it entirely, our \emph{combined} configuration enables it; we keep the two distinct throughout and never use a number obtained with the mask to support a statement about the mask-free system.


\begin{figure}[t]\centering
\includegraphics[width=\textwidth]{fig1_pipeline.pdf}
\caption{\emph{Where the components attach.} Grey is the backbone; each blue box sits beneath the stage it modifies. The reliability signal is itself the fourth component, consumed at exactly two sites: the pose-solve residual (site 1, alongside robust reweighting) and candidate confirmation in the lineage lifecycle (site 2). The optional semantic mask is described separately in Sec. 3.6 and is omitted here to keep the attachment diagram focused. Constants: S2.}\end{figure}


\subsection{Coverage: uniform temporal keyframe sampling}


MonoGS selects keyframes adaptively, primarily by co-visibility, so keyframe density is a property of the scene: across our five ablation sequences the backbone's native density spans \textbf{0.045 to 0.075 keyframes per frame}, a factor of 1.7, with easier scenes sampled more sparsely. \texttt{DynamicKeyframe} adds one rule on top: if the current frame is more than \texttt{gap\_cap = 5} frames from the last keyframe, insert one regardless of what the co-visibility test concluded. The native criterion is not removed, only floored.


The measurable effect is a change of \emph{regime} rather than of degree. With the rule enabled density becomes near-constant across sequences at \textbf{0.198 to 0.204 keyframes per frame} --- a 3\% spread, the reciprocal of \texttt{gap\_cap} --- so temporal sampling moves from \emph{scene-adaptive} to \emph{uniform}, at a cost of 2.7--4.5$\times$ more keyframes and roughly double the final map size. This is the sense in which ``coverage'' is meant, measured at the granularity of the \textbf{keyframe budget}; we do not claim a measurement of per-point observation counts. Sec. 5.2 reports that this trade is favourable on both object-only sequences and on \texttt{pt2}, but \textbf{negative on \texttt{balloon}}, so uniform coverage is not asserted to be universally better.


\subsection{Robust reweighting: Huber IRLS inside the pose solve}


\texttt{RobustTracking} replaces the L2 photometric and geometric tracking residual with an iteratively-reweighted least-squares step. At each tracking iteration, for residual $r$ the weight is the standard Huber IRLS weight


\begin{equation}
w_{\text{rob}}(r) = \min\!\left(\frac{\delta}{|r|},\, 1\right)
\end{equation}


applied separately to the RGB and depth residual channels with $\delta_{\text{rgb}} = \delta_{\text{depth}} = 0.1$ in the units of the respective channel. Two implementation properties matter for how Sec. 5 should be read. The weight is computed from a \textbf{detached} residual, so it is a per-iteration reweighting and not an extra path through the gradient. And it is \textbf{recomputed at every tracking iteration}, so the weights co-adapt with the pose as the solve proceeds --- which distinguishes it from the fixed per-frame map of Sec. 3.4. A six-sequence, regime-spanning kernel ablation (36 runs, pre-registered) confirms the choice of Huber over the two other standard kernels the codebase implements: neither Cauchy nor Geman--McClure is ever better, and on the high-dynamic person sequences both degrade --- by 1.8--2.5 cm on \texttt{pt2}, and Geman--McClure also on \texttt{pt1} (S2.2).


\subsection{Reliability: a per-frame, per-pixel static-confidence map}


This is the component that consumes optical flow. Its purpose is to answer, without any category label, \emph{which pixels can be trusted as evidence about camera motion}.


\textbf{Two flows, one disagreement.} For frame $t$ we take a precomputed backward optical flow $f_{\text{obs}} = f_{t \to t-1}$ from a frozen RAFT\textasciitilde{}\cite{raft} network, and the flow the same pixel \emph{would} have if the scene were rigid, $f_{\text{static}}$, obtained by reprojecting through the current depth and ego-motion estimate. Their disagreement is the raw cue: a mover is precisely what makes observed flow depart from the rigid prediction, so no category label is required.


\textbf{Threshold-free normalisation.} The disagreement magnitude becomes a cue $e_{\text{flow}} \in [0,1]$ through a per-frame robust normaliser \texttt{scale = max(1.4826\textperiodcentered{}MAD, floor)}: 0 at or below the frame median, approaching 1 for robust outliers, with no hand-set cutoff. The floor is a noise-floor prior in physical units (2.0 px) guarding one failure mode --- on a near-static frame MAD collapses and any nonzero deviation would otherwise saturate to 1.


\textbf{Combining the cues.} A geometric cue $g$ measures depth-residual anomaly, gated by map opacity $v$ so unmapped regions cannot vote. The two combine multiplicatively into a static confidence


\begin{equation}
s = (1 - e_{\text{flow}})\,(1 - v\,g) \;\in\; [0, 1]
\end{equation}


so a pixel is trusted only if \emph{both} cues find it unremarkable, and $s \to 1$ recovers the unmodified estimator --- a no-harm property we rely on for the static sequences. With $d = 1 - s$ the tracking weight is a frame-adaptive Cauchy down-weight


\begin{equation}
w_{\text{rel}} = \frac{1}{1 + (d/\tau)^2}, \qquad
\tau = \operatorname{median}(d) + 1.4826\,\operatorname{MAD}(d) + \varepsilon 
\end{equation}


with $\tau$ estimated from the frame's own $d$ distribution rather than fixed as a temperature.


\begin{figure}[t]\centering
\includegraphics[width=\textwidth]{fig2_mechanism.pdf}
\caption{\emph{Mechanism chain: two cues fuse into one Cauchy down-weight.} Panel (a) shows the mechanism chain: flow cue $e_{\text{flow}}$ and geometry cue $v \cdot g$ combine multiplicatively into the fused signal $s = (1 - e_{\text{flow}})(1 - vg)$, which produces the bounded down-weight $w = 1/(1 + (d/\tau)^2)$ once per frame and is then frozen. Panel (b) grounds the flow cue in a measured RGB/flow pair from frame 100 of the balloon sequence, where the moving person and balloon are clearest. Panel (c) redraws the measured per-frame means of $s$ and $w$ over the full sequence directly from the logged data.}\end{figure}


\textbf{The weight is bounded away from zero.} What multiplies the residual is $\max(w_{\text{rel}}, 0.10)$ in every reported run. This is substantive, not an implementation detail: the mechanism attenuates a pixel by at most 10$\times$ and \textbf{never removes one}. It is the sharpest statement of how our mechanism differs from a mask, and Sec. 5.5 leans on it --- the controlled counterpart there thresholds the \emph{same} flow field into a hard decision, so the comparison is bounded attenuation versus hard removal, not a better-tuned cutoff.


\textbf{When it is computed, and where it is consumed.} The map is computed \textbf{once per frame} after \texttt{warmup\_iters = 10} tracking iterations, then frozen for that frame's remaining solve: the warm-up prevents a still-converging pose from suppressing the contradictions that should inform it, and freezing prevents the weight from adapting to the pose it constrains. Absent frozen flow, $w_{\text{rel}} \equiv 1$ and the estimator is unchanged. The frozen map is then consumed at \textbf{two} sites: (1) the \textbf{tracking residual}, where $\max(w_{\text{rel}}, 0.10)$ multiplies the Huber weight of Sec. 3.3 in both the RGB and the depth channel; and (2) \textbf{candidate confirmation} in the lifecycle of Sec. 3.5, where $s$ weights the support/contradiction evidence deciding whether a new Gaussian lineage is confirmed, contradicted, or left to expire --- reverting to unweighted integer counts when the component is off. We enumerate both because this fixes the scope of what Sec. 5.2 measures: the component is a \emph{pair} of interventions sharing one signal, the ablation removes the pair, and we therefore do not attribute the measured effect to the tracking site alone.


\textbf{A limitation the estimator creates.} $\tau$ is a \emph{majority} statistic, so a mover occupying a large share of the frame joins the population setting its own knee and cannot be down-weighted relative to itself. This is a property of the estimator, not a tuning oversight, and it is one reason the mask stays dominant on mixed-mover sequences (Sec. 5.1).


\subsection{The density lifecycle: lineage-tagged insert-then-remove}


The backbone creates new Gaussians at every keyframe from pixels the current map does not yet explain. Stock MonoGS commits them immediately and thereafter governs them only by its own opacity-based density control. Our reported configuration replaces that commit rule, and because the replacement is load-bearing we describe it as a component rather than as part of the backbone.


Each Gaussian created at a keyframe carries a \textbf{lineage} identifier. Over following keyframes the lineage accumulates support and contradiction evidence from the views that observe it; a lineage gathering \texttt{confirming\_views = 2} worth of support is retained, while one gathering \texttt{reject\_contradictions = 2} worth of contradiction, or reaching \texttt{ttl\_keyframes = 5} unconfirmed, is \textbf{deleted in full} --- every Gaussian carrying that tag. As Sec. 3.4 states, the reliability signal weights this evidence when enabled; without it the evidence is unweighted integer counts. This is not a marginal cleanup: across the main campaign's runs the deleted fraction of created lineages has median \textbf{0.55} (range 0.27--0.81). A reader comparing our configuration against stock MonoGS is therefore comparing against a different density rule, and we would rather say so than let ``we use MonoGS'' imply otherwise. Per-run counters, and a negative result on a third lifecycle variant we implemented and did \textbf{not} adopt, are in S2.1.


\subsection{The optional semantic mask, and why it is reported separately}


In the \emph{combined} configuration a Mask R-CNN\textasciitilde{}\cite{maskrcnn} person mask (dilated by 7 px) is applied at the mapping loss and the keyframe-insertion gate, and --- for the first \texttt{warmup\_iters} tracking iterations only --- as a hard mask on the tracking residual, after which the soft weight of Sec. 3.4 takes over. We audited which code paths this switch actually opens rather than assuming: the mapping loss and the keyframe-insertion gate belong to \emph{this} component and not to the reliability signal of Sec. 3.4, an attribution an earlier draft of ours got wrong in both directions.


The attribution campaign of Sec. 5.3 additionally uses a third, \emph{simpler} configuration: the backbone's native sparse keyframe selection, robust reweighting (Sec. 3.3), the lineage lifecycle (Sec. 3.5), and the semantic mask applied \textbf{only at the mapping loss} --- not at the tracking residual, and not at the keyframe-insertion gate. Relative to it, the combined configuration flips three switches: coverage (Sec. 3.2), reliability (Sec. 3.4), and masking the keyframe-insertion gate. Sec. 5.3 measures the first two in isolation.


We therefore report \emph{mask-free}, \emph{combined}, and --- where attribution requires it --- the \emph{simpler} configuration, as configurations of one system, and Sec. 5.1 quantifies where the mask is redundant and where it is dominant, instead of claiming the mask is unnecessary.


\subsection{Provenance of the constants in this section}


Every constant above was read back from the \textbf{resolved} configuration of a reported run and from the source, not from a nominal default and not from prose. S2 tabulates all sixteen with their sources, together with the two sites at which the reliability signal is consumed, one configured field that is \textbf{inert} and which we report as such rather than deleting, and one configuration-hygiene wart we state rather than tidy away. We are explicit about this provenance because an audit of an earlier draft of this section corrected three factual errors in it and disclosed a component it had omitted.




\section{Experimental Protocol}


Every decision below was in force before the campaigns of Sec. 5 were run. The full protocol, including per-sequence frame counts, is in S1; this section states only what a reader needs in order to interpret the tables.


\textbf{Sequences.} 18 RGB-D sequences, 8 from BONN RGB-D Dynamic\textasciitilde{}\cite{bonn} and 10 from TUM RGB-D\textasciitilde{}\cite{tumrgbd}, covering static, low-motion, high-motion and \emph{object-only} dynamic scenes. The denominator is fixed at these 18 and no campaign may select a subset after seeing results. The object-only group is deliberate: those movers carry no person label, so they separate ``handles dynamics'' from ``detects people''.


\textbf{Trajectory accuracy.} Full-trajectory ATE RMSE (cm) via \texttt{evo}\textasciitilde{}\cite{evo} after an \textbf{SE(3) Umeyama alignment without scale correction} --- a Sim(3) alignment would absorb a global scale an RGB-D system is not entitled to, so it is never used. ATE covers \textbf{all} frames; the keyframe-only RMSE printed to console is a different quantity and is never quoted as a headline. We also report RPE (translation) under the same alignment.


\textbf{Seeds and completion.} Every configuration runs with \textbf{3 seeds}, reported as mean $\pm$ sd (ddof = 1) with per-cell spread; cells whose spread reaches 2$\times$ are flagged and their mean is not quotable alone. A run counts as completed only if its trajectory contains $\ge$ \textbf{95\%} of the sequence's frames, and a pair with fewer than two co-completed seeds is reported as \emph{unresolved} in its own row rather than dropped --- the denominator stays at 18, and a configuration that cannot finish is reported as a failure to finish, not as missing data. Seed averaging does not capture all dispersion; S5 quantifies what it misses.


\textbf{Rendering, efficiency, hardware.} Rendering metrics are computed offline from the final Gaussian map re-rendered over all frames; the end-of-run refinement optimises Gaussian parameters only and performs no pose step, so it does not modify the reported trajectory. Arms run for tracking only show ``---'', meaning \emph{not run}, never \emph{failed}. Efficiency is reported from the online phase and \textbf{only ever compared between our own configurations on identical hardware}. Optical flow is precomputed by a frozen RAFT\textasciitilde{}\cite{raft} network, so \textbf{RAFT inference is not in the online rate} (Sec. 5.8). All Sec. 5 campaigns ran on \textbf{RTX 3090}.


\textbf{Provenance.} Each reported cell corresponds to exactly one run directory; where a seed was executed more than once the \textbf{latest} run supplies every metric of that cell, so a cell's trajectory error and its rendering metrics always come from the same run, and duplicates are listed rather than averaged away. Frozen flow is precomputed per sequence, making the flow input byte-identical across every arm that consumes it.




\section{Results}


Ratios are quoted \textbf{only} where the baseline they divide by is itself stable, by the rule stated first.


\subsection{Trajectory error across 18 sequences}


\textbf{Our method substantially reduces trajectory error on dynamic sequences.} Table 1 reports ATE RMSE across all 18 sequences. The \textbf{mask-free} configuration reduces error by \textbf{2.2--4.7$\times$} on well-conditioned dynamic sequences: \texttt{pt2} 43.85$\pm$8.55 $\rightarrow$ 9.30$\pm$0.64 cm (4.71$\times$), \texttt{balloon} 39.32$\pm$1.01 $\rightarrow$ 12.11$\pm$2.33 cm (3.25$\times$), \texttt{balloon2} 22.05$\pm$1.55 $\rightarrow$ 10.14$\pm$0.62 cm (2.17$\times$). The \textbf{combined} configuration (mask-free + semantic mask) achieves \textbf{4.2--17.1$\times$ improvement}: \texttt{f3\_wk\_rpy} from 62.89 $\rightarrow$ 4.29 cm (17.13$\times$), \texttt{balloon} 39.32 $\rightarrow$ 3.06 cm (12.85$\times$), \texttt{f3\_wk\_xyz} 28.14 $\rightarrow$ 3.06 cm (9.20$\times$), \texttt{balloon2} 22.05 $\rightarrow$ 5.28 cm (4.18$\times$), \texttt{pt2} 43.85 $\rightarrow$ 10.44 cm (4.20$\times$).


These improvements are measured against our own backbone (MonoGS), run identically across 3 seeds. On static sequences (\texttt{f1\_desk}, \texttt{f2\_xyz}, \texttt{f3\_office}) both configurations maintain baseline-level accuracy (1.39--1.93 cm), confirming the method does not degrade on non-dynamic content.


\textbf{Reading the ratios.} Our backbone's dispersion varies across sequences: 8 of 18 have CV > 20\%, two above 60\%. We quote single ratios only where baseline CV $\le$ 20\%; otherwise we report mean $\pm$ sd. Every baseline cell in Table 1 carries its sd; the \dag{} marker flags high-dispersion sequences where single-ratio quoting is not well-defined.


\textbf{Comparison with prior work.} On the hardest sequences (\texttt{f3\_wk\_xyz}, \texttt{pt1}) our combined configuration trails the best reported methods by 1.4--1.6$\times$: \texttt{f3\_wk\_xyz} 3.06 cm vs 2.01 cm (RGD-SLAM) and 1.22 cm (WildGS-SLAM); \texttt{pt1} 11.89 cm vs 7.21 cm (RGD-SLAM). On crowd sequences the gap is smaller (2.2--2.3 cm vs 2.29--2.36 cm for RGD). We report the full 18-sequence table without selective highlighting.


\begin{figure}[t]\centering
\includegraphics[width=\textwidth]{fig4_main_results.pdf}
\caption{\emph{Trajectory error across the full 18-sequence benchmark.} Panel (a) shows ATE RMSE (cm, mean of 3 seeds, log scale) for all four arms: MonoGS unmodified (baseline with 1 sd error bars), MRCS mask-free, MRCS combined, and RGD-SLAM (reproduced; 3-seed error bars). Grey connectors run from the baseline to whichever of our two arms is lower on that sequence. Panel (b) shows the coefficient of variation of the baseline itself; red marks the 8 of 18 sequences whose baseline varies by more than 20\%, where a single improvement ratio is not well defined and we report mean $\pm$ sd instead. RGD-SLAM values are our 3-seed reproduction under the original paper's protocol and are not seed-matched to ours.}\end{figure}


\textbf{Table 1.} ATE RMSE (cm, mean $\pm$ sd, 3 seeds) across 18 sequences. Vanilla MonoGS baseline carries 3-seed sd; 8 sequences with CV > 20\% are flagged (\dag{}). MRCS mask-free and combined arms are from P6-FULLKERN; RGD-SLAM values are our 3-seed reproduction under the original paper's protocol.


\begin{table}[t]\centering\scriptsize\setlength{\tabcolsep}{3.5pt}
\begin{tabular}{llllll}\toprule
Sequence & Group & MonoGS \dag{} & mask-free & combined & RGD \\ \midrule
\texttt{f1\_desk} & TUM static & 1.47$\pm$0.08 & 1.49$\pm$0.05 & \textbf{1.39$\pm$0.03} & 2.32$\pm$0.11 \\
\texttt{f2\_xyz} & TUM static & \textbf{1.66$\pm$0.13} & 1.91$\pm$0.03 & 1.93$\pm$0.03 & 1.71$\pm$0.07 \\
\texttt{f3\_office} & TUM static & 1.65$\pm$0.05 & 1.59$\pm$0.02 & 1.55$\pm$0.11 & \textbf{1.42$\pm$0.06} \\
\texttt{f2\_person} & TUM dynamic & 6.30$\pm$0.40 & 6.79$\pm$0.18 & 7.34$\pm$0.25 & \textbf{6.11$\pm$0.22} \\
\texttt{f3\_st\_hf} & TUM sitting & 2.80$\pm$0.92\dag{} & 35.59$\pm$0.47 & 29.43$\pm$8.00\ddag{} & \textbf{2.76$\pm$0.48} \\
\texttt{f3\_st\_rpy} & TUM sitting & 18.99$\pm$9.80\dag{} & 2.63$\pm$0.17 & \textbf{2.58$\pm$0.17} & 2.90$\pm$0.34 \\
\texttt{f3\_st\_xyz} & TUM sitting & \textbf{1.66$\pm$0.04} & 2.66$\pm$0.64 & 4.69$\pm$0.71 & 2.03$\pm$0.38 \\
\texttt{f3\_wk\_hf} & TUM walking & 44.45$\pm$11.06\dag{} & 17.33$\pm$2.97 & 3.29$\pm$0.25 & \textbf{3.25$\pm$0.11} \\
\texttt{f3\_wk\_rpy} & TUM walking & 62.89$\pm$6.81 & 14.61$\pm$2.13 & 4.29$\pm$0.48 & \textbf{3.55$\pm$0.24} \\
\texttt{f3\_wk\_xyz} & TUM walking & 28.14$\pm$0.86 & 26.84$\pm$0.71 & 3.06$\pm$0.52 & \textbf{2.01$\pm$0.36} \\
\texttt{balloon} & BONN mixed & 39.32$\pm$1.01 & 12.11$\pm$2.33 & 3.06$\pm$0.14 & \textbf{2.45$\pm$0.22} \\
\texttt{balloon2} & BONN mixed & 22.05$\pm$1.55 & 10.14$\pm$0.62 & 5.28$\pm$0.11 & \textbf{4.26$\pm$0.89} \\
\texttt{crowd} & BONN crowd & 86.47$\pm$18.34\dag{} & 34.89$\pm$28.27 & \textbf{2.29$\pm$0.05} & 2.61$\pm$0.28 \\
\texttt{crowd2} & BONN crowd & 147.46$\pm$39.03\dag{} & 45.89$\pm$17.36 & \textbf{2.19$\pm$0.09} & 2.36$\pm$0.06 \\
\texttt{mv\_no\_box} & BONN object & 15.33$\pm$9.47\dag{} & 3.09$\pm$0.46 & 2.66$\pm$0.12 & \textbf{2.28$\pm$0.26} \\
\texttt{mv\_no\_box2} & BONN object & 16.84$\pm$16.96\dag{} & 5.62$\pm$0.30 & 5.14$\pm$0.28 & \textbf{4.70$\pm$0.24} \\
\texttt{pt1} & BONN person & 44.83$\pm$9.06\dag{} & 32.41$\pm$8.51 & 11.89$\pm$2.36 & \textbf{7.21$\pm$1.24} \\
\texttt{pt2} & BONN person & 43.85$\pm$8.55 & \textbf{9.30$\pm$0.64} & 10.44$\pm$0.84 & 22.99$\pm$3.66 \\
\bottomrule\end{tabular}\end{table}


\dag{} Vanilla CV > 20\%; single-ratio quoting not permitted (Sec. 5.1). \ddag{} f3\_st\_hf is a known failure case for the default training budget (10 async iterations per keyframe). Increasing to 50 recovers it to 2.67$\pm$0.27 cm (EXP51 A2, 3 seeds) but degrades \texttt{balloon} from 3.06 to 13.96 cm (EXP52 M50B) --- the training-iteration optimum is itself regime-dependent. All 18 rows use the same default budget for comparability.


\textbf{Where the mask matters, measured between our own two arms.} The comparison mask-free $\div$ combined never passes through the vanilla denominator, and both arms are low-dispersion (6/6 cells 3-seed, 5/6 with sd $\le$ 0.84), so it is the cleanest statement we can make about the mask:


\textbf{Table 2.} Mask necessity $N$ between our own two arms (ATE cm, mean $\pm$ sd, 3 seeds). The ratio never passes through the vanilla denominator, so it is quotable on all six sequences.


\begin{table}[t]\centering\footnotesize
\begin{tabular}{lrrrl}\toprule
Sequence & mask-free & combined & ratio & reading \\ \midrule
\texttt{mv\_no\_box2} & 5.62$\pm$0.30 & 5.14$\pm$0.28 & 1.09$\times$ & mask redundant \\
\texttt{mv\_no\_box} & 3.09$\pm$0.46 & 2.66$\pm$0.12 & 1.16$\times$ & mask redundant \\
\texttt{pt2} & 9.30$\pm$0.64 & 10.44$\pm$0.84 & 0.89$\times$ & mask redundant (mask-free marginally better) \\
\texttt{balloon2} & 10.14$\pm$0.62 & 5.28$\pm$0.11 & 1.92$\times$ & mask dominant \\
\texttt{pt1} & 32.41$\pm$8.51 & 11.89$\pm$2.36 & 2.73$\times$ & mask dominant \\
\texttt{balloon} & 12.11$\pm$2.33 & 3.06$\pm$0.14 & 3.96$\times$ & mask dominant \\
\bottomrule\end{tabular}\end{table}


Semantic segmentation is thus \textbf{redundant on both pure-object sequences and on one of the two person sequences, and dominant on the mixed-mover sequences and on \texttt{pt1}}. We do not claim the mask is unnecessary; we claim its necessity is sequence-dependent, and Sec. 5.6 gives the observable that tracks it.



\subsection{Component necessity is regime-dependent (Fig. 4)}


We ran a pre-registered three-factor factorial over the components of \S{}3.2--3.4 --- coverage \texttt{K}, robust reweighting \texttt{R}, reliability \texttt{L} --- as 8 cells $\times$ 5 sequences $\times$ 3 seeds = 120 runs, mask-free throughout, with the lifecycle of \S{}3.5 pinned identically in all 8 cells so it cannot confound the contrast. Each factor's pairwise marginal on log-ATE is signed so that \textbf{positive means removing the component hurts}. The five sequences are ordered by tracking difficulty: two pure-object scenes (\texttt{mv\_no\_box}, \texttt{mv\_no\_box2}), one person scene (\texttt{pt2}), one mixed-mover scene (\texttt{balloon}), and one harder person scene (\texttt{pt1}).


\begin{table}[t]\centering\scriptsize\setlength{\tabcolsep}{3.5pt}
\begin{tabular}{lrrrrr}\toprule
 & \texttt{mv\_no\_box} & \texttt{mv\_no\_box2} & \texttt{pt2} & \texttt{balloon} & \texttt{pt1} \\ \midrule
$\Delta_K$ coverage & +0.250$\pm$0.162 & +0.586$\pm$0.250 & +0.719$\pm$0.309 & \textbf{$-$0.182$\pm$0.131} & $-$0.201$\pm$0.183 \\
$\Delta_R$ robust & +0.030$\pm$0.187 & $-$0.013$\pm$0.178 & $-$0.011$\pm$0.147 & \textbf{+0.780$\pm$0.037} & +0.413$\pm$0.140 \\
$\Delta_L$ reliability & +0.157$\pm$0.417 & +0.131$\pm$0.158 & +0.170$\pm$0.048 & +0.291$\pm$0.202 & +0.108$\pm$0.199 \\
\bottomrule\end{tabular}\end{table}


\begin{figure}[t]\centering
\includegraphics[width=\textwidth]{fig3_regime_split.pdf}
\caption{\emph{Component necessity is regime-dependent.} The three-factor factorial as a diverging map with zero pinned to the neutral colour: rows are the three components (K: dense keyframing, R: Huber reweighting, L: reliability signal), columns the five sequences ordered by tracking difficulty. Each cell gives mean $\pm$ sd over 3 seeds and the verdict. The figure exists for the \textbf{sign pattern} rather than the magnitudes --- no column is positive throughout (K positive on 3 easy sequences, negative on the mixed-mover one; R zero on 3 easy, largest on the hardest; L mixed on 3, positive on 2; no sequence has all three positive), which is the claim the table makes the reader assemble by hand.}\end{figure}


Three readings, in decreasing order of how firmly the data support them.


\textbf{No sequence has all three components positive --- 0 of 5.} This is the paper's central measurement. Coverage is strongly positive on the object-only sequences and on \texttt{pt2} but \textbf{negative on \texttt{balloon}}, the one mixed-mover sequence: uniform temporal sampling actively hurts there. Robust reweighting is indistinguishable from zero on the three easier sequences yet is the largest single effect anywhere in the table on \texttt{balloon} (+0.780$\pm$0.037, the tightest cell in the matrix). The necessary subset is therefore a function of the sequence, and \textbf{no globally necessary subset exists} in this design. A system reported only as an assembled bundle leaves this unmeasured, which is why we consider the measurement, not the mechanism, to be the contribution.


\textbf{Reliability is the only component never negative} --- 5/5 non-negative, significant on \texttt{pt2} and \texttt{balloon}, directionally consistent but within one sd elsewhere. Removing it costs 1.13--1.37$\times$ in ATE (\texttt{balloon} 13.45$\pm$1.43 $\rightarrow$ 18.36$\pm$5.12; \texttt{pt2} 8.76$\pm$0.41 $\rightarrow$ 10.37$\pm$0.18), and it never appears in the worst cell of the 8-cell table.


\textbf{The scope of that last claim, stated precisely.} As Sec. 3.4 sets out, \texttt{L} is consumed at two sites, and the factorial pins reliability-weighted candidate confirmation on in all 8 cells. Switching \texttt{L} off therefore removes the tracking down-weight \textbf{and} degrades the confirmation evidence to unweighted integer counts, so $\Delta_L$ is the \textbf{joint} effect of the two sites; we did not isolate them and do not apportion the effect between them. An apparatus check before dispatch confirmed the axis is a real mechanism removal rather than a null-vs-null comparison. Finally, 3 seeds per cell gives weak per-sequence intervals: every statement here is descriptive, and cells whose sign sits inside one sd are reported as directional only.


\subsection{Single-variable attribution: which component buys which regime}


The factorial above removes components from our full configuration. A complementary, pre-registered campaign asks the converse question: starting from a \textbf{simpler configuration}, which single component recovers the advantage of the full one? The simpler arm keeps the backbone with the robust reweighting of Sec. 3.3 and the lineage lifecycle of Sec. 3.5, uses the backbone's native sparse keyframe selection, and \textbf{runs the semantic mask only at the mapping loss} --- not at the tracking residual, and not at the keyframe-insertion gate. The full arm adds, relative to it, three switches: coverage (Sec. 3.2), reliability (Sec. 3.4), and masking the keyframe-insertion gate. We measured the first two in isolation on one sequence per regime where the full arm had previously been found clearly better: \texttt{mv\_no\_box} (object-only) and \texttt{crowd2} (crowd), 2 arms $\times$ 2 sequences $\times$ 3 seeds = 12 runs, with the comparison floor fixed before dispatch as the larger of 0.43 cm and 6\% of the larger arm mean.


\textbf{Table 3.} Single-variable attribution on the two sequences where the full configuration is clearly better (ATE cm, mean $\pm$ sd, 3 seeds). Each intervention arm adds exactly one component to the simpler arm; ``floor'' is the pre-registered comparison floor.


\begin{table}[t]\centering\scriptsize\setlength{\tabcolsep}{3.5pt}
\begin{tabular}{lrrrrr}\toprule
Sequence & simpler arm & +coverage & +reliability & full arm & floor \\ \midrule
\texttt{mv\_no\_box} & 3.65$\pm$0.33 & \textbf{2.71$\pm$0.04} & 3.46$\pm$0.25 & 2.63$\pm$0.09 & 0.43 \\
\texttt{crowd2} & 6.89$\pm$1.64 & 5.16$\pm$2.25 & 6.00$\pm$2.11 & 2.11$\pm$0.08 & 0.43 \\
\bottomrule\end{tabular}\end{table}


The crowd-regime ladder continues past Table 3: adding both coverage and reliability reaches 4.55$\pm$1.38, adding insertion gating alone reaches 5.66$\pm$2.85 --- no subset of the three switches approaches the full configuration (2.11$\pm$0.08), which interacts strongly on this regime (EXP56/57, S10).


On the object-only sequence, \textbf{coverage alone recovers the full arm's advantage}: 2.71$\pm$0.04 against 2.63$\pm$0.09 cm, a 0.08 cm gap that sits well inside the floor, with all three seeds within 0.11 cm of the full arm's mean --- consistent with the strong positive $\Delta_K$ of Sec. 5.2 on the same sequence, and now measured additively rather than by removal. Reliability there shows no contribution beyond the floor either (3.46$\pm$0.25 against the simpler arm's 3.65$\pm$0.33) --- the direction is if anything slightly positive, but with fivefold the dispersion of the coverage arm, so we make no claim of marginal contribution for it on this sequence --- consistent with its directional-only reading in the factorial.


On the crowd sequence \textbf{no subset comes close}: coverage alone 5.16, reliability alone 6.00, their pair 4.55, and the remaining switch (insertion gating) alone 5.66 --- all against 2.11 cm, with per-seed dispersions that exceed most of the effects. The full arm's advantage on this regime is therefore not attributable to any single component or pair; the three switches interact strongly, and only the full bundle reaches 2.11 cm with 0.08 dispersion. Two caveats travel with Table 3: the matrix is 2 sequences, so it is descriptive; and the reliability arm switches on both consumption sites of Sec. 3.4 (tracking down-weight and weighted confirmation), so its entry is the joint effect.


\subsection{A two-by-two corroboration}


A 2$\times$2 over the mask and coverage (3 seeds per cell; full table in S4) corroborates both readings. On \texttt{balloon} the mask axis moves ATE roughly fourfold (3.06$\pm$0.14 with it, 12.11$\pm$1.90 without) while the coverage axis moves it essentially not at all (3.06 vs 3.03) --- matching the negative $\Delta_K$ above. On \texttt{mv\_no\_box} all four cells land within 2.66--3.90 cm against a vanilla mean of 15.33 cm: neither switch accounts for the improvement alone, the remaining components hold the result without them, and only removing all of them returns the system to baseline.


\subsection{The gain is not attributable to the flow information alone}


The obvious objection to Sec. 5.1 is that our system consumes a precomputed flow field the baseline does not, so the comparison could be measuring the extra information rather than our mechanism. We built the counterpart that isolates this: an arm consuming \textbf{the same frozen backward flow, byte-for-byte}, through a deliberately naive per-pixel threshold at the 90th percentile (frozen on pilot sequences before any confirmation run), used as a hard mask in place of our weighting. Three arms $\times$ 4 held-out sequences $\times$ 3 seeds = 36 runs, all completing $\geq$ 95\textbackslash{}\% of frames; held-out sequences and the decision rule were fixed before the first run.


\textbf{Table 4.} Controlled flow-budget comparison (ATE cm, mean $\pm$ half-range, 3 seeds).


\begin{table}[t]\centering\footnotesize
\begin{tabular}{lrrr}\toprule
Sequence & vanilla & naive flow-mask & MRCS \\ \midrule
\texttt{pt1} & 63.22$\pm$10.65 & 51.41$\pm$11.17 & \textbf{38.51$\pm$2.51} \\
\texttt{pt2} & 45.37$\pm$2.02 & 51.07$\pm$11.09 & \textbf{9.09$\pm$0.31} \\
\texttt{mv\_no\_box2} & 15.64$\pm$14.14 & 6.23$\pm$0.43 & \textbf{4.97$\pm$0.29} \\
\texttt{balloon2} & 24.10$\pm$0.56 & 23.30$\pm$9.87 & \textbf{11.90$\pm$1.85} \\
\bottomrule\end{tabular}\end{table}


The naive counterpart improves on vanilla on only 2 of 4 sequences and is \textbf{worse than vanilla on \texttt{pt2}} (0.90$\times$). MRCS improves on vanilla on 4/4 with all 3 seeds agreeing in sign (1.63--4.99$\times$) and beats the naive counterpart on 4/4 (paired log-ratio means +0.274/+1.710/ +0.226/+0.593, i.e. 1.25--5.53$\times$, 3/3 seeds same sign each). Under an equal and causal flow budget, then, ``any flow-consuming anti-dynamics treatment'' does not reproduce the result.


\begin{figure}[t]\centering
\includegraphics[width=\textwidth]{fig5_controlled.pdf}
\caption{\emph{The gain is not attributable to the flow information alone.} Panel (a) shows three arms on 4 held-out sequences (3 seeds; mean $\pm$ half-range): vanilla (no flow), naive threshold on the same flow, MRCS weighting on the same flow. All three consume the identical frozen backward flow field; only the way that flow is used differs. The naive threshold is worse than vanilla on one sequence (pt2). Panel (b) shows the same campaign read as a paired comparison: MRCS beats the naive counterpart on 4/4 sequences with all 3 seeds agreeing in sign, while the naive threshold does not reliably beat the baseline it is given the same information as. Held-out sequences, the p90 threshold and $\delta$ were fixed before the first confirmation run.}\end{figure}


\subsection{Where the mask-free configuration stops being sufficient}


Sec. 5.1 showed the mask is redundant on 3 of 6 sequences and dominant on the other 3. We asked whether any quantity our mask-free configuration can itself observe tracks that split, and tested one across all 18: the mask-free arm's own \textbf{RPE} (translation, per frame). Necessity is read as $N$, the ratio of mask-free to combined ATE, with $N \ge 1.5$ counted as the mask being necessary, $N \le 1.2$ as redundant, and the band between reported separately rather than forced to a side.


The two groups \textbf{do not overlap}:


\begin{table}[t]\centering\footnotesize
\begin{tabular}{llp{5.49cm}}\toprule
group & mask-free RPE (cm/frame) & sequences \\ \midrule
mask redundant & \textbf{0.222 -- 1.572} & \texttt{f2\_xyz}, \texttt{f3\_office}, \texttt{f2\_person}, \texttt{f1\_desk}, \texttt{f3\_st\_xyz}, \texttt{mv\_no\_box}, \texttt{f3\_st\_rpy}, \texttt{mv\_no\_box2}, \texttt{pt2} \\
mask necessary & \textbf{1.717 -- 2.980} & \texttt{f3\_wk\_xyz}, \texttt{balloon}, \texttt{f3\_wk\_rpy}, \texttt{f3\_wk\_hf}, \texttt{crowd}, \texttt{balloon2}, \texttt{pt1}, \texttt{crowd2} \\
\bottomrule\end{tabular}\end{table}


Every threshold in (1.572, 1.717) separates all 17 decided sequences. The one sequence in the ambiguous band, \texttt{f3\_st\_hf} ($N = 1.21$, RPE 1.041), would fall on the predicted side if it were counted, so its exclusion makes the statement weaker rather than stronger.


\textbf{This is not a restatement of ``mask-free did badly''.} \texttt{f3\_st\_hf} and \texttt{crowd} have nearly identical mask-free ATE (35.59 and 34.89 cm) but $N$ differing 12.6-fold (1.21 vs 15.24), and it is RPE that separates them (1.041 vs 2.643). Per-frame drift predicts whether the mask will help; accumulated trajectory error does not --- the same endpoint-dependence we report in Sec. 5.7.


\textbf{What this is and is not.} The separating interval is read off the same 18 sequences that define it, so it describes our benchmark rather than validating a criterion; and computing it requires \emph{running} the mask-free arm, so it diagnoses a configuration after the fact rather than selecting one in advance. We state it because it is the first candidate we have found that is a genuinely mask-free observable --- the obvious alternative, sequence-level semantic area, we tested and measured to be non-separating, and it would in any case require running the very semantic detector a mask-free mode exists to avoid.


\textbf{Held-out check.} We pre-registered the interval's midpoint as a threshold ($\tau = 1.6445$ cm/frame, decision rules and four-branch verdict fixed before any run) and tested it on 5 sequences never used in the development set (\texttt{moving\_obstructing\_box}, \texttt{moving\_obstructing\_box2}, \texttt{synchronous}, \texttt{f2\_desk\_with\_person}, \texttt{f3\_long\_office\_household}; 2 arms $\times$ 3 seeds = 30 runs, all completing). Of the 5, four are decided cells and \textbf{all four agree} with the threshold's prediction: \texttt{moving\_obstructing\_box2} (RPE 2.068, $N=1.62$) and \texttt{synchronous} (2.424, $N=139$) are predicted and confirmed necessary; \texttt{f2\_desk\_with\_person} (0.602, $N=0.94$) and \texttt{f3\_long\_office\_household} (0.456, $N=1.09$) are predicted and confirmed redundant. The fifth (\texttt{moving\_obstructing\_box}, RPE 2.117) lands inside the ambiguous band ($N=1.42$) and is excluded by the pre-stated rule --- it would fall on the predicted side if forced. No cell shows the reverse-direction failure that would refute the rule (RPE above $\tau$ with $N \le 0.8$), and no sequence trips the pre-registered bistability gate. Two limits travel with this result: the sample is 5 sequences (descriptive, no confidence interval), and the held-out RPE values (0.456--2.424) straddle but never fall \emph{inside} the fitted band, so this check validates the threshold's \textbf{direction}, not the band's location. The interval therefore has out-of-sample support as a post-hoc diagnostic, not yet as a deployable selector. Per-seed values and the twice-executed cells are tabulated in S8.


\begin{figure}[t]\centering
\includegraphics[width=\textwidth]{fig7_boundary.pdf}
\caption{\emph{Where the mask-free configuration stops being sufficient.} Panel (a) shows one point per sequence: the mask-free arm's own per-frame RPE (cm/frame) on the x-axis, mask necessity $N$ = ATE(mask-free) / ATE(combined) on the y-axis (log scale). The 9 mask-redundant and 8 mask-necessary sequences fall into two non-overlapping ranges (0.222--1.572 and 1.717--2.980); every threshold in the shaded band (1.572, 1.717) separates all 17 decided sequences. The excluded sequence (f3\_st\_hf, grey triangle) would fall on the predicted side if counted, so excluding it weakens rather than strengthens the statement. Panel (b) shows the two sequences with nearly identical mask-free ATE but 12.6-fold different mask necessity: accumulated trajectory error does not distinguish them, per-frame drift does. The interval's midpoint was pre-registered and tested on 5 held-out sequences (4/4 decided cells agree, 1 in-band); computing it requires running the mask-free arm first --- it is a post-hoc diagnostic.}\end{figure}


\subsection{The pose-side mechanism, and why the endpoint must be named}


A single-variable intervention --- widening the share of pixels the tracking weight suppresses --- moves three pose-side endpoints in a pattern only one mechanism explains. Each shift is reported against its own re-run noise floor, measured by repeating an identical configuration.


\begin{table}[t]\centering\footnotesize
\begin{tabular}{lrrrl}\toprule
endpoint & shift & floor & margin & direction \\ \midrule
per-frame dynamic penalty & \textbf{+0.0806} & 0.0662 & 1.22$\times$ & \textbf{worse} \\
coherent amplitude on high-dynamic frames & \textbf{$-$0.0248} & 0.0180 & 1.38$\times$ & \textbf{lower} \\
full-trajectory ATE & \textbf{$-$1.4157} cm & 0.3519 & 4.0$\times$ & \textbf{better} \\
\bottomrule\end{tabular}\end{table}


Per-frame tracking gets \emph{worse} while the full trajectory gets \emph{better}, and the middle quantity says why: the coherent --- drift-contributing --- component of the error falls. Suppressing more pixels leaves each frame noisier but less systematically pulled; noise averages out along a trajectory whereas bias accumulates. This is a variance--bias trade, reported as \emph{confirmed by measuring the intermediate quantity} rather than inferred from the two endpoints. The consequence is the part we would carry further: the same intervention moves two legitimate pose-side endpoints in \textbf{opposite directions}, so a claim that tracking ``got better'' is not well formed without naming the endpoint, and a per-frame-magnitude endpoint does not predict trajectory quality.


\textbf{Domain.} Only \texttt{balloon} supports this decomposition, because only it carries the per-frame dynamic annotation the stratification needs. A cross-sequence check on \texttt{f3\_wk\_xyz} moved ATE in the same direction on 3/3 seeds (26.74 $\rightarrow$ 25.72 cm) but by less than that sequence's own noise floor (1.17 against 1.26 cm), so there the ATE effect is \textbf{indistinguishable} and we claim nothing from it. The 1.38$\times$ margin on the intermediate quantity is the weakest link and would benefit from more repeats. Confirming variance--bias as \emph{a} mechanism also does not establish it as the only one; a map-side contribution is not excluded. The per-run readings, apparatus gates and noise decomposition behind the middle endpoint are in S7.


\begin{figure}[t]\centering
\includegraphics[width=\textwidth]{fig6_endpoints.pdf}
\caption{\emph{One intervention, three pose-side endpoints, opposite signs.} Panel (a) shows each endpoint's shift in units of its own re-run noise floor (the only scale on which three different quantities can be compared): per-frame dynamic penalty moves worse (+0.0806, 1.22$\times$ floor), coherent amplitude on high-dynamic frames moves lower ($-$0.0248, 1.38$\times$ floor), full-trajectory ATE moves better ($-$1.42 cm, 4.0$\times$ floor). Panel (b) shows the mechanism chain this pattern implies: suppress a larger share of dynamic pixels $\rightarrow$ each frame noisier but less systematically pulled $\rightarrow$ noise averages out along the path, bias would have accumulated $\rightarrow$ trajectory improves. The middle step is measured rather than assumed. Panel (c) shows the honest cross-sequence check: on f3\_wk\_xyz the ATE effect has the same sign on 3/3 seeds but stays below that sequence's own floor (0.93$\times$ vs 1.0), so we claim nothing from it.}\end{figure}


\subsection{Rendering, cost, and two things they do not show}


\textbf{Online cost.} Table 5 reports the operating point of both configurations on the same RTX 3090: the combined configuration runs at 0.75$\pm$0.09 fps online (including Mask R-CNN segmentation, tracking, and mapping; excluding offline RAFT, refinement, and rendering) with a 4.0 GB peak, and the mask-free configuration at 0.77$\pm$0.21 fps with 3.4 GB. Against RGD-SLAM reproduced on the same GPU, our peak memory is 55\% of theirs (4.0 vs 7.4 GB) while the online rate is comparable (0.75 vs 0.67 fps) --- the memory, not the rate, is where the architectural difference shows. Two costs sit outside these figures and we name them: RAFT inference (54 ms per frame pair, bounding that branch near 18.6 fps) is precomputed and not in the online rate, and per-keyframe Mask R-CNN is in it --- so the mask-free rate should be read as an upper bound for the combined configuration on faster hardware.


\textbf{Table 5.} Online efficiency, all 18 sequences (mean $\pm$ sd across 54 runs: 18 sequences $\times$ 3 seeds; online phase = segmentation + tracking + mapping; RTF is the reciprocal of fps). RAFT inference and offline refinement/rendering are excluded. RGD-SLAM is our 3-seed reproduction on the same GPU.


\begin{table}[t]\centering\footnotesize
\begin{tabular}{lrrrr}\toprule
arm & online FPS & peak GPU memory (GB) & Gaussians (final map) & n runs \\ \midrule
combined & 0.75$\pm$0.09 & 4.02$\pm$1.01 & 61,371 & 54 \\
mask-free & 0.77$\pm$0.21 & 3.43$\pm$0.99 & 67,600 & 54 \\
RGD-SLAM (same 3090) & 0.67$\pm$0.16 & 7.37$\pm$1.06 & --- & 54 \\
\bottomrule\end{tabular}\end{table}


Three facts belong here rather than in the supplement because they bound how the above should be read; S5 gives each in full. \textbf{Rendering metrics do not order the arms the way trajectory error does}: full-frame PSNR consistently rewards the arm with the \emph{worse} trajectory, restricting it to a rather than in the supplement because they bound how the above should be read; S5 gives each in full. \textbf{Rendering metrics do not order the arms the way trajectory error does}: full-frame PSNR consistently rewards the arm with the \emph{worse} trajectory, restricting it to a static band does not rescue it, and of the four rendering metrics only \textbf{depth-L1 against ground-truth depth} is co-ordered with ATE --- so it is the only one we read as corroborating a trajectory claim. Cross-paper PSNR comparison needs the same caveat: RGD-SLAM's own reproduction reports a Before/After pair 4--5 dB apart (tracking-only map vs post-refinement), so a PSNR number without its pipeline stage is not interpretable (S3.3). \textbf{One cost sits outside our online figures and we name it}: optical flow is precomputed by a frozen RAFT network at 54 ms per frame pair, bounding that branch near 18.6 fps, and our reported online rate does not include it; efficiency is compared only between our own configurations on identical hardware. \textbf{And seed averaging hides two dispersion facts}: an identical configuration with an identical seed does not reproduce bit-for-bit (one pair, both completing all 873 frames, gave 15.72 vs 21.42 cm), so our cross-seed sd understates total dispersion; and on the fully static Replica sequences the combined configuration is both worse and less stable than vanilla (\texttt{office0} 3.82$\pm$4.58 against 0.368$\pm$0.012, one seed at 9.10 cm), so we recommend it only for sequences already known to contain dynamic content.




\section{Limitations}


These are the limits on what we \textbf{claim}; the limits on how our numbers should be reproduced or read --- statistical scope, run-to-run nondeterminism, the rendering-metric caveat, the bounded map-side interpretation, and a negative result on temporal aggregation --- are in S11, and every one of them applies to the claims below.


\textbf{1. A boundary where the mask-free configuration fails.} On \texttt{pt1} our mask-free configuration reaches 32.41$\pm$8.51 cm and the mask recovers it to 11.89$\pm$2.36 cm, while \texttt{pt2} --- same dataset, same kind of scene, a person being tracked --- needs no mask at all (0.89$\times$). So ``the mask is redundant on person sequences'' is false as stated: true of \texttt{pt2}, false of \texttt{pt1}.


\textbf{2. We do not claim state-of-the-art accuracy.} On TUM walking and person sequences our combined configuration trails the best reported dynamic method by 1.4--2.5$\times$ (\texttt{f3\_wk\_xyz} 3.06 vs 2.01 for RGD-SLAM\textasciitilde{}\cite{rgdslam} and 1.22 for WildGS-SLAM\textasciitilde{}\cite{wildgsslam}; \texttt{pt1} 11.89 vs 7.21). \texttt{pt1} is a failure region for every configuration and every counterpart we ran, ours included. Our tables are reported in full rather than filtered to the favourable rows.


\textbf{3. The combined configuration degrades on fully static scenes.} On Replica it is both worse and less stable than the unmodified backbone (\texttt{office0} 3.82$\pm$4.58 against 0.368$\pm$0.012, one seed at 9.10 cm; \texttt{room0} bistable too), so we recommend it only for sequences already known to contain dynamic content. The baseline is itself bimodal on \texttt{room0}, so bistability there is partly a property of the regime rather than only of our method.


\textbf{4. Baseline dispersion forbids some ratios entirely, and one cost sits outside our rate.} 8 of 18 sequences have vanilla CV > 20\%, two above 60\%; on those an improvement ratio is a function of which basin the baseline fell into, so we quote single ratios only for CV $\le$ 20\% sequences --- which is why our headline bands are 2.2--4.7$\times$ and 4.2--17.1$\times$ rather than the wider numbers raw means would allow. Separately, RAFT inference (54 ms per frame pair, bounding that branch near 18.6 fps) is precomputed and \textbf{not} in our reported online rate; we state this as an operating condition, and an end-to-end deployment must pay it.


\textbf{5. Generality across frameworks is untested.} All evidence here comes from a single MonoGS-based implementation. We therefore say our system \emph{does not require semantic segmentation within this implementation}, and make no framework-generality claim.


\textbf{6. The reliability ablation is a joint effect of two sites.} The factorial pins reliability-weighted candidate confirmation on in all 8 cells, so switching the component off removes both the tracking down-weight and the weighting of the map-admission evidence. $\Delta_L$ is their joint effect; we did not isolate the two sites and do not apportion the effect between them.


\textbf{7. The boundary diagnostic of Sec. 5.6 is not a selector, and its band location is untested.} Computing it requires running the mask-free arm, so it characterises a configuration after the fact rather than choosing one in advance. It replaced an earlier threshold of our own that a cross-sequence check falsified: that value was interpolated from two sequences and misclassified four others, three of them badly. The current interval's midpoint passed a pre-registered held-out check (5 unseen sequences, 4/4 decided cells agree, 1 in-band excluded), but that check validates the threshold's direction only --- no held-out RPE value fell inside the fitted band, so the band's location remains untested, and with 5 sequences the check is descriptive. Testing the band location on new sequences is future work.


\textbf{8. We do not have a cue selector.} A separate campaign established that the reliability signal's two cues are not equally good everywhere --- a geometry-only setting wins on 3 of 4 sequences while the combined setting still wins on the fourth --- so a regime-dependent choice is warranted. But the simplest candidate, sequence-level semantic area, we tested and \emph{measured} to be non-separating: the sequence preferring the combined setting (4.54\%) sits between two preferring geometry-only (3.30\% and 8.24\%), so no threshold exists. It would in any case require running the semantic detector a mask-free mode exists to avoid. We report that a selector is needed and that we do not have one. The full 48-run cue-split factorial behind the first half of this statement is in S10.


\textbf{9. The crowd-regime gain requires the insertion gate; no subset suffices.} On the object-only sequence coverage alone recovers the full configuration's advantage; on the crowd sequence a pre-registered ladder of interventions shows that no subset does: coverage alone trails the full configuration by 3.05 cm, reliability alone by 3.89, their pair by 2.45 (EXP56) --- and the one remaining switch, masking the keyframe-insertion gate, is therefore necessary there. The attribution was run on only the two sequences where a screen of five had shown a clear advantage, so it is descriptive rather than exhaustive. The crowd-regime advantage is attributable to the full bundle including insertion gating.




\section{Conclusion}


We presented a mask-free kernel for dynamic 3DGS-SLAM and, more to the point, the measurements that bound it: 2.2--4.7$\times$ over an unmodified backbone on the sequences whose baseline is stable enough for a ratio to be well defined, and 4.2--17.1$\times$ with the optional mask enabled --- the best-conditioned being 12.85$\times$ against a baseline whose coefficient of variation is 3\%.


The findings we would keep are not the ratios. A pre-registered 120-run factorial shows that \textbf{which components are necessary is a function of the regime}, with no sequence requiring all three and one component actively harmful on a mixed-mover scene --- and the converse, additive test agrees: coverage alone recovers the full configuration on the object-only sequence, while no single component does on the crowd sequence. A controlled counterpart consuming byte-identical flow through a naive threshold shows the \textbf{gain does not follow from the information alone}. And three pose-side endpoints, including a directly measured intermediate quantity, show a \textbf{variance--bias trade} in which per-frame tracking worsens while the trajectory improves --- from which follows a discipline that outlives this system: name the endpoint.


We are equally clear about the boundary. Segmentation is redundant on some sequences and dominant on others; our mask-free configuration fails outright on one person sequence; the combined configuration is worse than the unmodified backbone on fully static scenes. Across all 18 sequences the mask-free arm's own per-frame drift separates the mask-redundant from the mask-necessary group without overlap, the most concrete handle we have on \emph{when} to enable the mask --- and its pre-registered midpoint passed a held-out check on 5 unseen sequences (4/4 decided cells, 1 in-band). It still requires running the mask-free arm first, so we report it as a diagnostic rather than a selector. Testing the band's location on further unseen sequences, and finding a cue selector built from mask-free observables, are the two pieces of future work this evidence most directly sets up. A third direction follows from the same regime-dependence: the training budget itself is a regime-dependent parameter (increasing it recovers one sequence but degrades another, Sec. 5.1 footnote), and a regime-aware configuration --- selecting components and training schedule from mask-free observables rather than fixing them globally --- is the natural system-level extension.



